\origpage{162}{ROBERTSON AND SEYMOUR}
\noindent $T \setminus f_i$ containing $u_i$ ($i = 1, 2$). Then $V(T_0) = V(T_1) \cup V(T_2) \cup \{u\}$ and $V(T_1) \cap V(T_2) = \emptyset$. Now the result holds for $(u_1, f_1)$ and $(u_2, f_2)$ by the inductive hypothesis. Moreover, since $E - \alpha(u_i, f_i) = \alpha(u, f_i)$ ($i = 1, 2$) and $(T, \alpha)$ is exact, it follows that
\begin{align*}
(E - \alpha(u_1, f_1)) \cup (E - \alpha(u_2, f_2)) &= E - \alpha(u, f) \\
(E - \alpha(u_1, f_1)) \cap (E - \alpha(u_2, f_2)) &= \emptyset.
\end{align*}
The result follows. $\square$

\stmt
(3.4) \textit{If there is no bias extending $\mathscr{A}$ then there is an exact tree-labelling over $\mathscr{A}$.}

\textit{Proof.} By (3.2), it suffices to prove that there is a tree-labelling over $\mathscr{A}$. Suppose that $E = \emptyset$. If $\emptyset$ is efficient, let $T$ be a two-vertex tree, and let $\alpha(v, e) = \emptyset$ for both incidences $(v, e)$ of $T$; $(T, \alpha)$ is the required tree-labelling. If $\emptyset$ is not efficient, then $\mathscr{A}$ is a bias, a contradiction. Thus we may assume that $E \ne \emptyset$. Choose $x \in E$, and let $\mathscr{B}$ be the set of all efficient sets $B \subseteq E$ with $x \notin B$; then $\mathscr{B}$ is a bias. Since $\mathscr{B}$ does not extend $\mathscr{A}$, it follows that $\mathscr{A} \ne \emptyset$.

We proceed by induction on the number $N$ of efficient sets $X \subseteq E$ such that neither $X$ nor $E - X$ is a subset of any member of $\mathscr{A}$. We suppose first that $N = 0$. Let $\mathscr{B}$ be the set of all efficient sets which are subsets of members of $\mathscr{A}$. Since $\mathscr{A} \subseteq \mathscr{B}$, $\mathscr{B}$ is not a bias. But for every efficient set $X$, either $X \in \mathscr{B}$ or $E - X \in \mathscr{B}$ since $N = 0$. Thus there exist $X_1, X_2, X_3 \in \mathscr{B}$ with $X_1 \cup X_2 \cup X_3 = E$. Let $T$ be the tree with four vertices $t_0$, $t_1$, $t_2$, $t_3$ and edges $e_i$ with ends $t_0$, $t_i$ ($i = 1, 2, 3$). Define $\alpha(t_0, e_i) = X_i$, $\alpha(t_i, e_i) = E - X_i$ ($i = 1, 2, 3$). Then $(T, \alpha)$ is a tree-labelling over $\mathscr{A}$, as required.

Thus we may assume $N > 0$. Choose an efficient set $X \subseteq E$ such that neither $X$ nor $E - X$ is a subset of any member of $\mathscr{A}$, and subject to that with $X$ minimal. Since $\mathscr{A} \ne \emptyset$, $X \ne \emptyset$. Let $\mathscr{A}_1 = \mathscr{A} \cup \{X\}$, $\mathscr{A}_2 = \mathscr{A} \cup \{E - X\}$. Since there is no bias extending $\mathscr{A}$, there is no bias extending $\mathscr{A}_1$ or $\mathscr{A}_2$. From our inductive hypothesis there are exact tree-labellings $(T_1, \alpha_1)$ over $\mathscr{A}_1$ and $(T_2, \alpha_2)$ over $\mathscr{A}_2$. A leaf $t$ of $T_1$ is bad if $\alpha_1(t, e) \ne E$ and $\alpha_1(t, e) \cup A \ne E$ for all $A \in \mathscr{A}$, where $(t, e)$ is an incidence, and we define the bad leaves of $T_2$ similarly. Now if $t$ is a bad leaf of $T_1$ and $(t, e)$ is an incidence, then $\alpha_1(t, e) \cup X = E$ and so $E - \alpha_1(t, e) \subseteq X$. If $E - \alpha_1(t, e) \ne X$, then from our choice of $X$, either $E - \alpha_1(t, e) \subseteq A$ for some $A \in \mathscr{A}$ or $\alpha_1(t, e) \subseteq A$ for some $A \in \mathscr{A}$. In the first case $\alpha_1(t, e) \cup A = E$, a contradiction, since $t$ is bad. In the second case $E - X \subseteq \alpha_1(t, e) \subseteq A$, contrary to our choice of $X$. Thus $E - \alpha_1(t, e) = X$, for every bad leaf $t$. Since $X \ne \emptyset$, it follows from (3.3) that there is at most one bad leaf in $T_1$.
